\chapter{Background and Related Work}

The delivered game sits in a recognisable design space between a small-scale tactical skirmish game and a card-driven command system. However, the project was never intended to compete with content-heavy strategy RPGs or large collectible card games. From the start, the team treated related work as a guide for framing interaction patterns rather than a source of gameplay scope inflation. That position is consistent with both the coursework brief and the early one-pager, which emphasised a readable, teenager-friendly tactical prototype rather than a broad commercial feature set \cite{cw1Specification,onePager}.

\section{Design Context}

The core design context was a lightweight tactical grid battle in which the player controls a very small team and spends a limited hand of cards each turn. This places \projectname{} closer to a command-oriented tactics prototype than to a deckbuilding game. Positioning, range, targeting, and action ordering are the main decision variables, while content breadth is intentionally constrained. The project therefore uses cards as a readable action language, not as a collectible system with large pools, rarity, or synergy chains \cite{onePager}.

Unity was selected because the game depends on both world-space interaction and screen-space feedback. The board, tile highlights, unit placement, and combat feedback all live in the same runtime environment as the HUD, menu flow, and result presentation. That makes Unity a better fit than a document-style prototype or a purely web-first interface for this coursework. It also allowed the team to demonstrate software-engineering concerns that matter for game development in practice, such as separating turn rules from view state, isolating assets, and validating behaviour across both scene logic and UI flow.

Just as importantly, the team chose to keep the product as one polished encounter rather than stretching toward multiple maps, progression systems, or a larger content pipeline. This was not simply a feature cut; it was a design decision to preserve clarity and testability. A compact battle loop gave the team enough room to implement and validate selection order, card resolution, enemy response, battle-end transitions, and usability refinements without diluting the engineering evidence with unfinished systems \cite{cw1Specification}.

\section{Reference Strategy}

The team used external references selectively and at two different levels. At the gameplay level, the project borrowed only broad genre patterns: grid occupancy, ranged versus melee positioning, and card-labelled actions. No external gameplay framework was imported into the runtime. The battle loop, flow coordination, grid logic, enemy response, and HUD behaviour were authored within the project codebase under the main runtime scripts, which keeps ownership boundaries clear and makes the resulting software easier to explain in the report.

At the asset level, the team deliberately adopted a `CC0-first` strategy during the post-WR3 visual pass so that the final package would remain easy to attribute and redistribute. The imported first-slice UI theme uses Kenney packs for panels, buttons, and icons, while later visual work fell back to team-authored presentation code for unit silhouettes and board dressing when a coherent local character pack was not yet ready for stable import \cite{postWr3OpenAssetRegister}. This distinction matters: the UI art is third-party and documented, but the presentation wiring, layout logic, and gameplay-facing behaviour remain original coursework implementation.

The main consequence of this reference strategy is that related work shaped form more than function. External material helped the team avoid placeholder visuals and supported a clearer tabletop-fantasy identity, but it did not replace the need for local engineering decisions. The delivered system is therefore best understood as an original tactical vertical slice built with carefully bounded reference use, not as an assembly of imported mechanics.

\section{Related-Work Boundary}

Because this is a coursework prototype rather than an academic game-studies paper, the related-work discussion remains intentionally practical. The most relevant comparison points are genre conventions and asset-source choices, not a long survey of published games. The report therefore focuses on the references that materially affected implementation: the initial design brief, the one-pager that constrained scope, and the open-licensed asset sources that informed the final presentation direction \cite{cw1Specification,onePager,postWr3OpenAssetRegister}.
